\documentclass[12pt, twoside]{article}
%\documentclass[12pt, twoside]{article}
\usepackage[letterpaper, margin=1in, headsep=0.2in]{geometry}
\setlength{\headheight}{0.6in}
%\usepackage[english]{babel}
\usepackage[utf8]{inputenc}
\usepackage{microtype}
\usepackage{amsmath}
\usepackage{amssymb}
%\usepackage{amsfonts}
\usepackage[nomessages]{fp} %\FPeval{\var-name}{2*sin(pi/6)}
\usepackage{siunitx} %units in math. eg 20\milli\meter
\usepackage{yhmath} % for arcs, overparenth command
\usepackage{tikz} %graphics
\usetikzlibrary{quotes, angles, arrows, arrows.meta}
\usepackage{pgfplots}
\usepgfplotslibrary{fillbetween}
\pgfplotsset{compat=1.16}
\usepackage{graphicx} %consider setting \graphicspath{{images/}}
\usepackage{parskip} %no paragraph indent
\usepackage{enumitem}
\usepackage{multicol}
\usepackage{venndiagram}

\usepackage{fancyhdr}
\pagestyle{fancy}
\fancyhf{}
\renewcommand{\headrulewidth}{0pt} % disable the underline of the header
\raggedbottom
\hfuzz=2mm %suppresses overfull box warnings

\usepackage{hyperref}

\fancyhead[LE]{\thepage}
\fancyhead[RO]{Markscheme}
\fancyhead[LO]{La Scuola d'Italia / Dr.\ Huson / IB Math 11 \\* 28 October 2026}

\begin{document}
\subsubsection*{2.6 Early Finishers: Lines, Exponents, Radicals and Sequences --- Markscheme}

\emph{Extension set, scored separately from the quiz. Calculator allowed.
Version (a) repeats problems 3 and 4 of quiz version (b), with the same
marking.}

\textbf{\large Version (a)} \hfill \emph{Total: 19 marks.}

\begin{enumerate}[itemsep=0.3cm]

%--- Problem 1: Same-base equation; integers a, b from a surd product; equation quadratic in 3^x [11 marks] ---
\item[1.] [Maximum mark: 11]

\textbf{(a)} $\left(5^{2}\right)^{x + 1} = \left(5^{3}\right)^{2x - 3}$ \hfill \textbf{M1}

$2x + 2 = 6x - 9$ \hfill \textbf{A1}

$x = \dfrac{11}{4}$ \hfill \textbf{A1}

\emph{Accept $2.75$. M1 for writing both sides as powers of $5$. Wrong
methods: $2x + 1 = 6x - 9$ or $2x + 2 = 6x - 3$ (a bracket not multiplied
out) gives $\frac{5}{2}$ or $\frac{5}{4}$: M1A0A0.}

\textbf{(b)} $2a + 2b\sqrt{3} + a\sqrt{3} + 3b = 12 + 7\sqrt{3}$ \hfill \textbf{M1}

$2a + 3b = 12$, \quad $a + 2b = 7$ \hfill \textbf{A1}

$a = 3$, \quad $b = 2$ \hfill \textbf{A1}

\emph{Accept the conjugate route:
$\frac{12 + 7\sqrt{3}}{2 + \sqrt{3}} \times \frac{2 - \sqrt{3}}{2 - \sqrt{3}}
= \frac{24 - 12\sqrt{3} + 14\sqrt{3} - 21}{1} = 3 + 2\sqrt{3}$ (M1A1A1).
Accept trial with the expansion of $(2 + \sqrt{3})(3 + 2\sqrt{3})$ shown as
a check (M1A1A1). A bare answer earns M0A0A1.}

\textbf{(c)} $9^{x} = \left(3^{x}\right)^{2} = u^{2}$, so $3u^{2} - 82u + 27 = 0$ \hfill \textbf{M1}

$(3u - 1)(u - 27) = 0$

$u = \dfrac{1}{3}$ \hfill \textbf{A1}

$u = 27$ \hfill \textbf{A1}

$3^{x} = \dfrac{1}{3} \Rightarrow x = -1$ \hfill \textbf{A1}

$3^{x} = 27 \Rightarrow x = 3$ \hfill \textbf{A1}

\emph{M1 for the quadratic in $u$. The $u$ values may come from factorizing,
the formula or the GDC. Stopping at $u = \frac{1}{3}, 27$ earns
M1A1A1A0A0. Wrong method: reading $3^{x}$ as $3x$ gives $x = \frac{1}{9}, 9$:
A0A0 for those values.}

\newpage

%--- Problem 2: Arithmetic vs geometric: expressions, least n with v_n > u_n, least n with geometric sum > arithmetic sum [8 marks] ---
\item[2.] [Maximum mark: 8]

\textbf{(a)} $u_{n} = 7 + 3(n - 1)$ \hfill \textbf{A1}

$v_{n} = 2(1.5)^{n - 1}$ \hfill \textbf{A1}

\emph{Accept $u_{n} = 3n + 4$.}

\textbf{(b)} $2(1.5)^{n - 1} > 7 + 3(n - 1)$ \hfill \textbf{M1}

$u_{7} = 25$, $v_{7} = 22.8$; \quad $u_{8} = 28$, $v_{8} = 34.2$ \hfill \textbf{A1}

$n = 8$ \hfill \textbf{A1}

\emph{M1 for the inequality, or a GDC table or graph of both sequences. A1
for the values on both sides of the change. FT from (a): $v_{n} = 2(1.5)^{n}$
gives $n = 7$.}

\textbf{(c)} $\dfrac{2\left(1.5^{n} - 1\right)}{1.5 - 1} > \dfrac{n}{2}\left(2(7) + 3(n - 1)\right)$ \hfill \textbf{M1}

$n = 9$: $149.8 < 171$; \quad $n = 10$: $226.7 > 205$ \hfill \textbf{A1}

$n = 10$ \hfill \textbf{A1}

\emph{M1 for both sum formulas (or a GDC table of both sums). A1 for the
sums on both sides of the change. Wrong method: comparing terms instead of
sums gives $n = 8$: M0A0A0. FT from (a).}

\end{enumerate}

\newpage

\textbf{\large Version (b)} \hfill \emph{Total: 22 marks.}

\begin{enumerate}[itemsep=0.3cm]

%--- Problem 1: Lines with a parameter k: parallel (two values), perpendicular, intersection when one line is vertical [10 marks] ---
\item[1.] [Maximum mark: 10]

\textbf{(a)} gradients $-\dfrac{k}{2}$ and $-\dfrac{3}{k - 1}$ \hfill \textbf{A1}

$-\dfrac{k}{2} = -\dfrac{3}{k - 1}$ \hfill \textbf{M1}

$k(k - 1) = 6$, \quad $k^{2} - k - 6 = 0$, \quad $(k - 3)(k + 2) = 0$

$k = 3$ \hfill \textbf{A1}

$k = -2$ \hfill \textbf{A1}

\emph{Accept comparing coefficients, $\frac{k}{3} = \frac{2}{k - 1}$, for the
A1M1. Finding $k = 3$ by inspection only, with no equation, earns
A0M0A1A0.}

\textbf{(b)} $\left(-\dfrac{k}{2}\right)\left(-\dfrac{3}{k - 1}\right) = -1$ \hfill \textbf{M1}

$3k = -2(k - 1)$ \hfill \textbf{A1}

$k = \dfrac{2}{5}$ \hfill \textbf{A1}

\emph{Accept $0.4$. Wrong method: product of the gradients $= +1$ gives
$k = -2$: M0A0A0.}

\textbf{(c)} $L_{1}$: $x + 2y = 7$, \quad $L_{2}$: $3x = 5$ \hfill \textbf{M1}

$x = \dfrac{5}{3}$ \hfill \textbf{A1}

$y = \dfrac{8}{3}$ \hfill \textbf{A1}

\emph{M1 for substituting $k = 1$ into both equations ($L_{2}$ is vertical).
Accept $1.67$ and $2.67$.}

\newpage

%--- Problem 2: Geometric sequence with terms in k: quadratic for k with surd roots, r in surd form, sum to infinity [12 marks] ---
\item[2.] [Maximum mark: 12]

\textbf{(a)} $\dfrac{k + 1}{k - 3} = \dfrac{3k - 1}{k + 1}$ \hfill \textbf{M1}

$k^{2} + 2k + 1 = 3k^{2} - 10k + 3$ \hfill \textbf{A1}

$2k^{2} - 12k + 2 = 0$, \quad $k^{2} - 6k + 1 = 0$ \hfill \textbf{AG}

\emph{M1 for equating the two ratios, or $(k + 1)^{2} = (k - 3)(3k - 1)$.
A1 for both sides expanded correctly.}

\textbf{(b)} $k = \dfrac{6 \pm \sqrt{36 - 4}}{2} = \dfrac{6 \pm 4\sqrt{2}}{2}$ \hfill \textbf{M1}

$k = 3 + 2\sqrt{2}$ \ or \ $k = 3 - 2\sqrt{2}$ \hfill \textbf{A1}

\emph{A1 for both values in exact form. $5.83$ and $0.172$ earn M1A0.}

\textbf{(c)} $r = \dfrac{k + 1}{k - 3} = \dfrac{4 + 2\sqrt{2}}{2\sqrt{2}}$ \hfill \textbf{M1A1}

$= \dfrac{\left(4 + 2\sqrt{2}\right)\sqrt{2}}{4} = \dfrac{4\sqrt{2} + 4}{4}$ \hfill \textbf{M1}

$r = 1 + \sqrt{2}$ \hfill \textbf{A1}

\emph{First M1 for the ratio of consecutive terms, either pair. Second M1
for rationalizing. Wrong method: the inverted ratio
$\frac{k - 3}{k + 1}$ gives $\sqrt{2} - 1$: M0A0M1A0. FT from (b).}

\textbf{(d)} $r = \dfrac{4 - 2\sqrt{2}}{-2\sqrt{2}} = 1 - \sqrt{2}$ \hfill \textbf{A1}

$|r| = 0.414\ldots < 1$, so the sum to infinity exists \hfill \textbf{R1}

$S_{\infty} = \dfrac{-2\sqrt{2}}{1 - \left(1 - \sqrt{2}\right)} = \dfrac{-2\sqrt{2}}{\sqrt{2}}$ \hfill \textbf{M1}

$S_{\infty} = -2$ \hfill \textbf{A1}

\emph{R1 for $|r| < 1$ (or $-1 < r < 1$) with their value of $r$. Wrong
method: the inverted ratio gives $r = -1 - \sqrt{2}$ (no sum to infinity)
or uses $\sqrt{2} - 1$ to get $2 + 2\sqrt{2}$: A0, R1 only if their
$|r| < 1$, M1A0. FT from (b).}

\end{enumerate}
\end{document}
